\subsection{Mailman-style shared-sum scoring}
\label{sec:8-mailman}

JHQ's finite sign patterns suggest reusing partial sums across
candidates, following the Mailman idea \cite{mailman}. A four-bit half-code
has only 16 patterns, so many candidates can share intermediate work.
The test freezes the primary codes and candidate streams on Vogue,
arXiv and OpenAI-3072. Atomic buckets and private histograms at group
widths 4, 6 and 8 are compared with packed two-table scoring, using
128 probes and batches of 8 and 1,000.

The best replacement reaches only 0.187--0.581 of the scorer baseline's
throughput. Profiles show concentrated bucket updates, while private
histograms introduce more state and shared-memory updates. Arithmetic
reuse is insufficient when obtaining it costs more than the saved
lookup work. These are scorer measurements, not an integrated search
comparison. If $f$ is the baseline time fraction spent scoring and $s$
the replacement's relative scorer speed, the unchanged-rest model gives
\[
 \frac{T'}T=1+f\left(\frac1s-1\right)>1
 \quad(0<f\leq1,\ 0<s<1).
\]
Thus even this model predicts no whole-search gain. Integration stopped;
other finite-alphabet implementations remain possible. The left panel
of \cref{fig:8-negative} shows the six scorer comparisons.

\par\medskip\noindent\begin{minipage}{\linewidth}
\centering\captionsetup{hypcap=false}
\includegraphics[width=\linewidth]{ch8/negative-evidence.pdf}
\captionof{figure}[Mailman and rank-oriented coding]{Mailman scorer throughput (left) and offline shortlist size (right).}
\label{fig:8-negative}
\end{minipage}\par\medskip

\subsection{Rank-oriented coding experiments}
\label{sec:8-rank}

Minimising reconstruction error need not preserve neighbour order.
FP64 offline replay tests weighted scores, median and list-dependent
thresholds, learned losses and linear minimum mean-square-error decoding,
with 400 calibration, 200 validation and 400 test queries. On arXiv,
median thresholds reduce the smallest tested shortlist reaching recall
0.97775 from 875 to 375 candidates; the other two workloads fail their
respective targets. These are the offline counts in the right panel of
\cref{fig:8-negative}.

A separate GPU follow-up changes the primary representation and retrains
the residual codes. At its own arXiv target of 0.9766, throughput improves
$1.40\times$ at batch 8 and $1.21\times$ at batch 1,000. Vogue instead
peaks at recall 0.9630, below its 0.9644 target, and OpenAI-3072 nearly
ties. The offline count reduction and GPU speedup use different queries,
states and targets; they do not establish a threshold-only speedup.
Changing primary codes also changes what the residuals must reconstruct.

\subsection{Removing a shared mean}
\label{sec:8-centring}

A global data mean can consume sign-code resolution. Rotation preserves
that mean's norm: for orthogonal $\Pi$,
\[
 \bar y=\frac1N\sum_i\Pi x_i=\Pi\bar x,
 \qquad \|\bar y\|^2=\bar x^\top\Pi^\top\Pi\bar x=\|\bar x\|^2.
\]
In real arithmetic, subtracting the same mean $\mu$ from queries and
database vectors preserves distances,
$\|(q-\mu)-(x-\mu)\|^2=\|q-x\|^2$,
but changes the quantisation problem.

A follow-up compares the archived original index with a fresh index
trained on centred data, alternating both in one session. On arXiv,
whose mean norm is 0.837, full-pool calibration at 128 probes changes
$\alpha=200$ to $\alpha=64$, and the recall-0.97 frontier improves
$1.266\times$. At 512 probes the selected budget becomes 100, not 64.
The change is not uniformly beneficial: Vogue's high-budget recall at
512 probes falls from 0.9918 to 0.9896, and OpenAI-1536's selected budget
rises from 8 to 16 despite a nearly unchanged recall-0.97 throughput
ratio of 0.995. The arXiv result supports mean removal as a useful
representation change. Retraining also changes centroids and residuals,
so repeated, equally retrained controls are needed to isolate its
contribution. These runs extend the analysis, not the main comparison
with other systems.

\subsection{A safe bound can still prune too little}
\label{sec:8-residual-bound}

A separate FP64 pilot constructs a lower bound on squared distance to
the reconstructed vector using only primary codes. For each sign
pattern, a box encloses residual coordinates from all indexed vectors.
Shifting it by the primary reconstruction gives a lower bound via
distance from the query. Across seven search
settings and 72.8 million candidate checks, no violation is observed.
Yet median bound-to-distance ratios are only 0.13--0.24. Even an ideal
final threshold leaves almost all probed candidates for refinement,
usually far more than the existing shortlist. The pilot therefore
stopped before an integrated GPU implementation. Safety alone is
insufficient: pruning must be selective enough to repay evaluating the
bound and reading surviving residuals. A GPU version would also need
rounding-safe envelopes.

\subsection{Further alternatives and stopping decisions}
\label{sec:8-stopped}

\Cref{tab:8-stopped} summarises representative earlier alternatives.
They use their own versions and controls rather than the final
cross-system protocol. GPU timings and offline cost estimates are
identified separately; the latter are not additional system-frontier
measurements.

\par\medskip\noindent\begin{minipage}{\linewidth}
\centering\captionsetup{hypcap=false}
\captionof{table}[Alternative designs and stopping decisions]{Alternative designs and stopping decisions.}
\label{tab:8-stopped}
\footnotesize\setlength{\tabcolsep}{4pt}
\begin{tabularx}{\linewidth}{@{}>{\raggedright\arraybackslash}p{.24\linewidth}>{\raggedright\arraybackslash}p{.24\linewidth}X@{}}
\toprule
Alternative & Expected saving & Evidence and decision \\
\midrule
Larger selection buffer & Fewer compaction sorts & GPU tests lose up to 24\% throughput; added shared memory reduces the possible resident blocks in affected cases \\
Fused scan and refinement & Fewer launches and intermediate transfers & GPU tests lose up to 15\%; the saved transfer is small relative to the combined kernel's resource demand \\
Partial-score early exit & Skip remaining subspaces & GPU changes range from $-2.8$ to $+4.2\%$; warp-wide checks give no consistent gain \\
Table-free sign scoring & Remove lookup-table storage & GPU throughput falls 30--52\% despite the smaller shared-memory footprint \\
Progressive primary reads & Read full codes only for survivors & Offline replay and traffic estimates expose scattered second-stage reads; no new GPU search kernel was timed \\
Local-first refinement & Avoid a global primary merge & Offline quota policies do not beat global selection at equal work; a favourable oracle allocation is unavailable online \\
\bottomrule
\end{tabularx}
\end{minipage}\par\medskip
